% ================================================================
\section{Discussion}
\label{sec:discussion}
% ================================================================

We organise the discussion around the three guiding questions, then state what the evidence says about mechanisms and about evaluation practice. All statements refer to one architecture (a noise-conditioned DnCNN inside PnP-PGD), two training distributions on $28\times28$ digits and one CT training run; the Limitations section lists the limits.

\paragraph{(1) Does the learned prior keep its advantage over TV when the image distribution changes?}
Not for far shifts. On $28\times28$ digits the in-distribution advantage is about $3.4$\,dB over box-constrained TV ($+3.8$\,dB over TV-ADMM; Sections~\ref{sec:e1e2} and~\ref{sec:seeds-results}). It halves for a near shift (EMNIST letters, $+1.9$\,dB over TV-ADMM) and reverses for an MNIST-trained prior on Fashion-MNIST ($-3.1\pm0.3$\,dB over TV-ADMM, three seeds) and on the Shepp--Logan phantom ($-4.3\pm0.1$\,dB), with the same signs against box-constrained TV (Sections~\ref{sec:e4} and~\ref{sec:seeds-results}). The reversal persists under Poisson noise ($-2.0$ to $-3.1$\,dB and $-4.4$ to $-6.1$\,dB; Section~\ref{sec:dose-results}), although in SSIM it appears on the phantom but not on Fashion-MNIST.
A prior trained on Fashion-MNIST beats TV on Fashion-MNIST ($+2.2$\,dB) and, with the same tuned hyperparameters, on MNIST ($+1.4$\,dB) and EMNIST letters ($+0.4$\,dB), and is close to TV on the phantom ($-0.2$\,dB). What matters is therefore how the training distribution relates to the test distribution; that the Fashion-trained prior transfers better is observed, but its cause (for example a greater variety of grey levels and shapes) is untested.
On real CT slices a CT-trained prior tested on CT has no consistent accuracy advantage over TV-long to begin with. On non-anatomical phantoms its detection advantage over TV-long is significant in one of six conditions, against three of the six corresponding CT conditions, and its SSIM deficit against TV-long grows from $0.004$--$0.027$ to $0.095$--$0.190$, while its PSNR deficit stays below $0.75$\,dB (Section~\ref{sec:e19}).

\paragraph{(2) Can we find lesion-like structure that the ground truth does not support?}
Yes, in some settings. For the MNIST-trained prior on Fashion-MNIST it produces $0.20$--$0.66$ false lesions per image, more than TV on seven of eight lesion types even against the box-constrained TV (Section~\ref{sec:lesion-results}); under Poisson noise the same holds for lesions with $\sigma_l\ge0.75$\,px in 11 of 12 dose and lesion-type cells, against the unconstrained TV only (Section~\ref{sec:lesion-poisson}). On non-anatomical phantoms at quarter dose it adds a few false blobs per phantom relative to FBP-Hann (significant at $K=120$ at both response thresholds, and at $K=30$ only for $\gamma=1.5$) and, at $K=120$, more than TV (Section~\ref{sec:e19}).
In distribution the picture is different. On digits with Gaussian noise it produces at most $0.01$ false lesions per image ($0.06$--$0.07$ for some lesion types at $30$ photons per ray) but attenuates inserted lesions. On real slices it produces $0.24$--$1.38$ false blobs per slice, fewer than FBP-Hann in the four conditions where this was tabulated, and the strongest blobs lie mostly where FBP-Hann and TV also produce them (Section~\ref{sec:e18}).
The false-blob count of TV is itself unstable: it varies $26$-fold across the phantom conditions and for the smallest lesions exceeds that of the learned prior (Sections~\ref{sec:lesion-results},~\ref{sec:lesion-poisson},~\ref{sec:e19}), so ``more than TV'' is a statement about a specific PSNR-tuned TV solver.
We therefore do not claim that the prior creates false structure in general, nor that it is safe: its behaviour moves with noise level, lesion scale and condition, and 40 patients with one noise draw per condition cannot bound rare events.

\paragraph{(3) Do the conclusions hold across noise levels and conditions?}
Only in part. The in-distribution advantage is largest at moderate noise ($+3.7$\,dB at $100$--$300$ photons per ray) and shrinks at both ends of the dose range (Section~\ref{sec:dose-results}). At $\sigma\le0.01$ the digit experiments do not determine the ranking, which changes with the data model by up to $5.5$\,dB (Section~\ref{sec:crime-results}), so we make no claim there. Under angular mismatch the learned prior stays ahead throughout the tested range, but its margin falls from $+3.86$ to $+2.98$\,dB for $4^\circ$ jitter and to $+1.49$\,dB for a $4^\circ$ offset; the latter is mostly a rigid rotation of the reconstruction that a pixel-wise metric penalises (Section~\ref{sec:e3}), so we do not read it as evidence that the learned prior is the less robust method.
On CT the signs of the PSNR, SSIM and false-blob differences against TV-long change with dose and number of angles without a monotone trend (Sections~\ref{sec:e16},~\ref{sec:e18}). The sign of the detection advantage over TV-long is the same in all nine conditions (significant in six), but it is specific to the 300-iteration, PSNR-tuned TV: against the 150-iteration TV it is negligible at full dose for $K=30$ and $60$ ($+0.006$ and $0.000$; Section~\ref{sec:e17}).

\paragraph{Accuracy depends on the operating regime and on the baseline.}
A single accuracy number for a learned reconstruction method is uninformative without its regime, data model and baseline. The headline gap on digits moves by about $0.5$\,dB with the choice between unconstrained and box-constrained TV and, at low noise, changes sign with the data model. On CT it depends on the iteration count given to TV and on tuning details: the apparent deficit of PnP at $K=60$ fell from $0.85$ to $0.17$\,dB at full dose when the tuning set grew from 6 to 20 slices and the starting image was changed (Section~\ref{sec:history3}).
On clinical-size slices the advantage found on digits disappears: PnP-const differs from TV-long by $-0.93$ to $+0.33$\,dB, with a lower SSIM in eight of nine conditions (significantly in seven). Both iterative methods improve on FBP-Hann ($0.9$--$1.4$\,dB for PnP-const, $0.8$--$2.3$\,dB for TV-long), so the PSNR gain over FBP-Hann is not specific to the learned prior (Section~\ref{sec:e16}). The best PnP variant, started from the TV solution, is highest in PSNR in seven of nine conditions by at most $0.49$\,dB, but the variants were introduced after the first PnP-const result, and two of its tuned settings lie on a grid boundary, one of them with almost no denoising ($\sigma_d=0.0039$ at $K=120$, full dose; Table~\ref{tab:params3}; Sections~\ref{sec:e15},~\ref{sec:e16}). TV-long and the PnP variants are not matched in cost or in tuned degrees of freedom.
We did not test why the digit and CT results differ. Candidate explanations are the more varied anatomy, which makes TV a weaker model than for digits, and a small denoiser trained once for $15\,000$ steps on $8\,960$ of $35\,820$ slices (a 17-layer variant did not help at $K=60$), so that the CT result cannot separate the prior from this network.

\paragraph{The iteration count acts as a regularisation parameter.}
In PnP-PGD with an unconstrained denoiser the accuracy peaks after a handful of iterations and then degrades, while the data residual stays at the noise level ($0.994$ to $1.019$ times its expected value between the selected iterate and $k=400$; Section~\ref{sec:e5}). Of the error growth, $74\%$ lies in directions with $0.1\le s_i/s_1<0.3$. Without the denoiser, the gradient step with $\eta=1/L$ would contract the distance to the data-fitting solution in direction $i$ by $(1-(s_i/s_1)^2)^k$, which is slow in these directions. The selected $k$ therefore determines how much of them is left to the denoiser, and it is a regularisation hyperparameter, not an implementation detail.
That the degradation is semi-convergence, in which further iterations fit noise in poorly determined directions, is our interpretation and not a tested mechanism: the selected $k$ is also $5$ at $\sigma=0$ and we did not record the iteration curve on noise-free data. The behaviour is consistent with the observation that convergence guarantees for PnP need constraints on the denoiser~\citep{ryu2019pnp}; we did not measure the Lipschitz constant of ours. In distribution the converged iterate is still better than TV on MNIST, whereas under shift its error grows to $3.7$--$4.8$ times that of TV (Section~\ref{sec:nullspace-results}).

\paragraph{The reversal under shift is a property of the prior, but not of the null space alone.}
The matched-prior control isolates the denoiser: in all runs the tuned hyperparameters are identical and only the weights differ (Section~\ref{sec:seeds-results}). The effect exceeds the seed-to-seed variation ($\le0.3$\,dB) by an order of magnitude and persists under four data models. It is not removed by re-tuning $(\sigma_d,k)$ within our grids: on Fashion-MNIST none of $45$ PnP settings reaches TV's PSNR, for one denoiser (Section~\ref{sec:oppoint-results}).
The error decomposition adds that the reversal is not simply null-space invention: $93\%$ (Fashion-MNIST) and $89\%$ (phantom) of PnP's excess error over TV lies outside the weakly determined subspace, and a signature of a prior that fills that subspace with up to three times the true energy, misaligned with the truth, appears only on the phantom (Section~\ref{sec:nullspace-results}). ``Hallucination'' in the null-space sense is therefore an incomplete description of the failure.

\paragraph{PSNR is not task performance, in either direction.}
In distribution on digits the PSNR-tuned PnP reconstruction is better than TV by about $4$\,dB but detects small inserted lesions less reliably (paired AUC $0.04$--$0.23$ below TV, also below the box-constrained TV), because it attenuates them (Section~\ref{sec:lesion-results}). Re-tuning trades attenuation against false lesions: for the MNIST-trained prior no operating point is competitive with TV on PSNR, detection and false lesions at once, on MNIST or on Fashion-MNIST. The Fashion-trained prior does leave such a point on both Fashion-MNIST and MNIST (about $+1.0$\,dB, a paired AUC $0.02$--$0.06$ below TV and few false lesions). We hypothesise that this reflects the variety of its training distribution and not whether a test set is formally in distribution; with one denoiser per prior this is untested (Section~\ref{sec:oppoint-results}).
On CT the sign is reversed: PnP detects better than TV-long in all nine conditions yet no better than FBP-Hann, and the PSNR-tuned TV-long, $0.8$--$2.3$\,dB better than FBP-Hann in PSNR, detects worse than it in all nine (significantly in seven) (Section~\ref{sec:e17}). A comparison reporting only PSNR or SSIM would have ranked the methods differently from this detection statistic, whichever method it favours.
We did not test why the sign of the detection difference differs between digits and CT; the lesion scale in pixels, the noise texture and the smoothing of the PSNR-tuned operating points are candidates. The CT operating points were tuned for PSNR, the lesions are small and of high contrast, and the observer knows their location, so the conclusions for detection-tuned operating points, low-contrast lesions, other observers and search tasks are untested.

\paragraph{Evaluation protocol changes the answer.}
Besides the choice of TV and its iteration count, three protocol choices affected results by amounts comparable to or larger than some of the differences we interpret. The data model (inverse crime or finer grid) changes the low-noise gap by up to $5.5$\,dB. The seed range of the Fashion-MNIST gap ($-2.83$ to $-3.43$\,dB) is of the size of the $0.5$\,dB box-constraint effect, although small against the shift effect itself. And choosing the lesion operating points on evaluation images would flatter re-tuning: the best of the $45$ settings with PSNR at least TV's, chosen on the test set, reaches a paired AUC of $0.898$, against $0.849$ for the point chosen on held-out validation images (Table~\ref{tab:oppoint}).
The shift sensitivity we observe is of the kind reported for other learned reconstruction methods~\citep{darestani2021shift}; those studies concern other methods and we make no comparison. Replacing the dense operator by a matrix-free projector with an exact adjoint changes no reconstruction at $28\times28$ by more than $3\times10^{-4}$ in relative norm, and replacing Gaussian by Poisson noise leaves the qualitative picture of the digit experiments unchanged (Sections~\ref{sec:mf-results},~\ref{sec:dose-results}); the Poisson-weighted data term showed no consistent benefit, so we claim only that the unweighted scheme is adequate in this setting.
We suggest, as practices and not as tested conclusions, that studies of learned priors report: a baseline with the same constraints and an iteration count tuned (or compute matched) like the learned method's, hyperparameters tuned on held-out data for every method, several training seeds, at least one mismatched and one matched prior, a task-based metric next to PSNR and SSIM, and an explicit data model.